% Part: model-theory
% Chapter: models-of-arithmetic
% Section: computable-models

\documentclass[../../../include/open-logic-section]{subfiles}

\begin{document}

\olfileid{mod}{mar}{cmp}
\section{Model Aritmetika Komputabel}

\begin{explain}
Model standar~$\Struct{N}$
nduweni rong sipat kang becik.
Ranahe yaiku wilangan asli~$\Nat$,
tegese unsur-unsuré pancen
prekara kang arep dirembug
nganggo basa aritmetika,
lan term angka standar~$\num{n}$
pancen nunjuk~$n$.
Sipat becik liyane:
interpretasi kabeh simbol
nonlogis saka~$\Lang{L_A}$
iku \emph{komputabel}.
Fungsi penerus, panjumlahan,
lan pingan kang dadi
$\Assign{\prime}{N}$,
$\Assign{+}{N}$,
lan $\Assign{\times}{N}$
iku fungsi wilangan
kang komputabel.
(Upamane bisa diitung
dening mesin Turing,
utawa ditegesi nganggo
rekursi primitif.)
Relasi kurang-saka
ing~$\Struct{N}$,
yaiku $\Assign{<}{N}$,
uga bisa diputusake.

Model nonstandar saka
teori aritmetika kayata
$\Th{Q}$ lan
$\Th{PA}$ kudu
ngemot unsur nonstandar.
Mula ranahe lumrahe
ngemot !!{element}s
saliyane~$\Nat$.
Nanging saben
!!{structure} kang
bisa didhaptar lan ranahe tanpa wates
bisa dibangun ing
sembarang himpunan kang
!!{denumerable},
kalebu~$\Nat$. Cathetan cakupan SOL6-F025: struktur winates ora bisa diangkut kanthi isomorfisme menyang ranah kaetung tanpa wates. Model Q lan PA ing andharan iki nduweni ranah tanpa wates, mula syarat kasebut katindakake.
Mula uga ana model
nonstandar kanthi
ranah~$\Nat$.
Ing model~$\Struct{M}$
kaya mangkono,
mesthi paling ora
sawenehing wilangan
ora nindakake peran
kang lumrah,
amarga sawenehing~$k$
kudu beda saka
$\Value{\num{n}}{M}$
kanggo saben~$n \in \Nat$.
\end{explain}

\begin{defn}
!!^a{structure}~$\Struct{M}$
kanggo $\Lang{L_A}$
iku \emph{komputabel}
yen lan mung yen
$\Domain{M} = \Nat$
lan $\Assign{\prime}{M}$,
$\Assign{+}{M}$,
$\Assign{\times}{M}$
iku fungsi komputabel,
dene $\Assign{<}{M}$
iku relasi kang
bisa diputusake.
\end{defn}

\begin{ex}\ollabel{ex:comp-model-q}
Elinga struktur
$\Struct{K}$
saka
\olref[mdq]{ex:model-K-of-Q}.
Ranahe yaiku
$\Domain{K} = \Nat
\cup \{a\}$
lan interpretasine
\begin{align*}
  \Assign{\Obj{0}}{K} & = 0\\
  \Assign{\prime}{K}(x) & =
  \begin{cases}
    x+1 & \text{yen $x\in \Nat$}\\
    a & \text{yen $x = a$}
  \end{cases}\\
  \Assign{+}{K}(x, y) & =
  \begin{cases}
    x+y & \text{yen $x$, $y \in\Nat$}\\
    a & \text{yen ora mangkono}
  \end{cases}\\
  \Assign{\times}{K}(x, y) & =
  \begin{cases}
    xy & \text{yen $x$, $y \in\Nat$}\\
    0 & \text{yen $x=0$ utawa $y=0$}\\
    a & \text{yen ora mangkono}\\
  \end{cases}\\
  \Assign{<}{K} & =
  \Setabs{\tuple{x,y}}{x, y \in \Nat \text{ lan } x<y} \cup
  \Setabs{\tuple{x,a}}{x \in \Domain{K}}
\end{align*}
Cathetan koreksi sumber OLPL-108:
peubah kang diiket ing
himpunan pasangan kapindho
yaiku x, padha karo
komponen kapisan pasangane.
Nanging $\Domain{K}$
iku !!{denumerable}
lan mula ekwinumerik
karo~$\Nat$.
Upamane,
$g\colon \Nat \to \Domain{K}$
kanthi $g(0) =
a$ lan $g(n) = n-1$
kanggo $n>0$
iku !!a{bijection}.
Cathetan koreksi sumber OLPL-109:
rumus n-1 nggawa
wilangan positif menyang
kabeh wilangan asli,
dene rumus n+1 ing sumber
ngilangake rong nilai awal.
Kita bisa nggunakake
pemetaan iki kanggo
ngasilake isomorfisme
antarane model anyar~$\Struct{K'}$
saka~$\Th{Q}$
lan $\Struct{K}$.
Ing $\Struct{K'}$,
kita kudu menehi fungsi
lan relasi kang beda
marang simbol-simbol
saka~$\Lang{L_A}$,
amarga !!{element}s
beda saka~$\Nat$
nglakoni peran
wilangan standar
lan nonstandar.

Mligine, $0$
saiki nglakoni
peran~$a$,
dudu wilangan standar
paling cilik.
Wilangan standar
paling cilik saiki~$1$.
Mula kita tetepna
$\Assign{\Obj{0}}{K'} = 1$.
Fungsi penerus
uga saiki beda:
yen diwenehi
wilangan standar,
yaiku $n > 0$,
asile isih
$n+1$.
Nanging $0$
saiki nglakoni
peran~$a$,
kang dadi penerus
kanggo awake dhewe.
Mula
$\Assign{\prime}{K'}(0) = 0$.
Kanggo panjumlahan
lan pingan,
kita uga nduweni
\begin{align*}
\Assign{+}{K'}(x, y) & =
  \begin{cases}
    x+y-1 & \text{yen $x$, $y >0$}\\
    0 & \text{yen ora mangkono}
  \end{cases}\\
  \Assign{\times}{K'}(x, y) & =
  \begin{cases}
    1 & \text{yen $x = 1$ utawa $y = 1$}\\
    xy - x - y + 2 & \text{yen $x$, $y > 1$}\\
    0 & \text{yen ora mangkono}\\
  \end{cases}
\end{align*}
Lan
$\tuple{x, y} \in \Assign{<}{K'}$
yen lan mung yen
$x < y$
lan $x > 0$
lan $y > 0$,
utawa yen
$y = 0$.

Kabeh fungsi iki
komputabel minangka
fungsi wilangan asli,
lan $\Assign{<}{K'}$
iku relasi kang
bisa diputusake
ing~$\Nat$.
Nanging fungsi-fungsi iki
beda saka fungsi
penerus, panjumlahan,
lan pingan lumrah
ing~$\Nat$,
lan $\Assign{<}{K'}$
uga beda saka
relasi~$<$
ing~$\Nat$.
\end{ex}

\begin{prob}
Wenehana !!a{structure}~$\Struct{L'}$
kanthi $\Domain{L'} = \Nat$
kang isomorfis
karo~$\Struct{L}$
saka
\olref[mod][mar][mdq]{ex:model-L-of-Q}.
\end{prob}

\begin{explain}
\Olref{ex:comp-model-q}
nuduhake yen $\Th{Q}$
nduweni model nonstandar
komputabel kanthi
ranah~$\Nat$.
Nanging asil ing ngisor
iki nuduhake yen
kahanan mangkono
ora dumadi kanggo
model saka~$\Th{PA}$
(lan mula uga
model saka~$\Th{TA}$).
\end{explain}

\begin{thm}[Teorema Tennenbaum]
$\Struct{N}$
iku siji-sijine model
komputabel saka~$\Th{PA}$
nganti isomorfisme.
Cathetan koreksi sumber OLPL-110:
kaunikan iki miturut
isomorfisme; salinan
komputabel kang mung
beda jeneng unsur
uga bisa ana.
\end{thm}
  
\end{document}
